\begin{textsample}{POS Dim 1 – vem\_ed\_17 – Score 5.00 – vem\_ed\_17/t079.txt}  \label{ex:f1_pos_136}
BOAS PRÁTICAS Organização, \textbf{flexibilidade} e engajamento estão entre as palavras-chave para o sucesso de um Projeto Colaborativo Internacional (PCI / Cesu).
Veterana em Intercâmbios Virtuais, a professora de Espanhol Lilian de Souza (Fatecs Americana, Bragança Paulista e Itu) desenvolve PCIs desde 2018.
Sua primeira experiência \textbf{foi} com a Universidad de Monterrey (UDEM / México).
Na ocasião, “um dos professores envolvidos se dispôs a fazer um site para alocar \textbf{todas} as informações do projeto e eu adorei a ideia”, relembra Lilian. “Posteriormente, tomei a iniciativa de estruturar um site numa plataforma gratuita para fazer um teste com os projetos de \textbf{que} participava.
Tanto os alunos como os professores gostaram e sinalizaram como um caminho fácil para encontrar as orientações sobre os PCIs”, comenta.
Com os sites, Lilian organiza as informações básicas de \textbf{cada} projeto (etapas, datas limites, orientações para entregas das \textbf{tarefas}, links de reuniões \textbf{síncronas}). \textbf{É} assim \textbf{que} Lilian consegue conduzir quatro PCIs simultaneamente no primeiro semestre de 2023: 1) Fatec Americana e Universidad Cundinamarca (Colômbia) – PCI “Sobre Culturas” mescla estudantes dos cursos de Produção Têxtil (Fatec Americana) e Psicologia (Universidad Cundinamarca).
As \textbf{tarefas} em equipe enfocam aspectos culturais de ambos os países, além de proporcionar aos estudantes brasileiros a prática de língua espanhola de maneira efetiva com falantes nativos – os alunos colombianos de Miguel Nicholls.
Site do PCI: sites.google.com/view/brasilcolombia/sobre-culturas

% matched lemmas: cada, flexibilidade, que, ser, síncrono, tarefa, todo
\end{textsample}
